\section{限制、连接与未来：dataset attention 的代价是什么}

NPT 的表达力来自让 datapoints 彼此可见，同一设计也制造最大系统瓶颈。若 batch 中有 $b$ 个 datapoints、每行 flatten 后 hidden dimension 为 $h=de$，ABD attention map 大小为 $b\times b$，主要计算可近似写为
\[
\operatorname{Cost}_{\mathrm{ABD}}=O(b^2h),
\qquad
\operatorname{Memory}_{\mathrm{attn}}=O(b^2).
\]

\begin{itemize}
  \item $b$：实际参与一次 relational lookup 的 datapoint 数；
  \item $h$：每个 row 的 hidden width；
  \item 增大 $b$ 提升候选 context 覆盖，却以 quadratic cost 增长；
  \item 随机 mini-batch 降低成本，但会遗漏 batch 外的潜在关键 neighbors。
\end{itemize}

\lecturefigure{slide-28-summary-future-work.jpg}{NPT 总结：实验、quadratic scaling 限制与 future work。}{Stanford Online 官方视频 01:00:42--01:05:36。}

\begin{knowledgebox}{读图：结论与开放问题分开看}
左侧是本讲已经展示的结果：dataset-as-input、datapoint self-attention、tabular/image benchmarks 与 intervention evidence。右侧是尚未解决的限制：scaling，以及 multi-task、few-shot、semi-supervised、domain adaptation 等新任务。Future work 不是当前能力清单。
\end{knowledgebox}

\subsection{与 GNN 的连接}

把每个 datapoint 看作 node，ABD 相当于在 fully connected graph 上做 learned message passing。区别在于许多 GNN 预先给定 sparse edges，而 NPT 从 attention score 中软学习关系；但现代 graph models 也可以学习 edges，因此二者不是绝对分界。

\teachervoice{Q\&A 中讲者接受 fully connected graph 的类比，并指出 NPT 可以被理解为 discovering relational structure。Neil Band 提到 Neural Relational Inference：当 edges 未知时，attention/message passing 都在学习潜在交互。}

\begin{warningbox}{Fully connected graph 不等于关系已被正确发现}
允许任意两点通信只是 hypothesis space。模型仍可能学到 shortcut、spurious similarity 或近似 uniform attention；需要 corruption、intervention、OOD evaluation 与可扩展 sparse retrieval 才能判断关系是否有用。
\end{warningbox}

\subsection{Small data 与 large data 的不同价值}

Small data 的优势是能把全数据集放进一次 forward，模型获得近乎完整 context；large data 则迫使我们引入 retrieval、sparsity、kernel approximation 或 learned representatives。后者虽然损失精确全连接 attention，却可能把 NPT 连接到现代 retrieval-augmented learning：参数负责学习检索与变换规则，外部 datapoints 负责保存可更新事实。

\teachervoice{讲者在最后提出一个尖锐问题：大型 parametric models 有多少参数其实在“存数据”？若显式 lookup 能承担部分记忆，参数或许可以更专注于学习怎样检索、组合与推理。这个问题后来在 retrieval-augmented systems 中变得更加重要。}

\begin{knowledgebox}{Scaling 路线图}
\begin{tabularx}{\textwidth}{L{0.24\textwidth}>{\raggedright\arraybackslash}X >{\raggedright\arraybackslash}X}
\toprule
方向 & 收益 & 新风险 \\
\midrule
Random mini-batch & 实现简单、训练稳定 & 关键 datapoint 可能不在 batch，context 随机。 \\
Sparse / local attention & 将 $b^2$ 降为少量 edges & 需要定义或学习 candidate graph。 \\
ANN / learned retrieval & 先检索 top-$k$ rows，再做 NPT reasoning & Retriever recall 成为上游瓶颈。 \\
Representative points & 用 prototypes / inducing points 压缩数据集 & 压缩可能丢失罕见样本与局部细节。 \\
Kernel / low-rank approximation & 近似全局 attention & 近似误差与硬件效率需共同评估。 \\
\bottomrule
\end{tabularx}
\end{knowledgebox}

\subsection{从固定数据集走向 meta-learning}

本讲实验为 fixed-dataset supervised learning，但 input contract 天然适合 context-based extension：把 support set、unlabeled set 与 query set 放入统一矩阵，通过 masks 指定任务。如果跨 datasets 训练，NPT 可能学习通用的 within-dataset algorithm；这与 in-context learning、few-shot learning、semi-supervised learning 和 domain adaptation 形成直接连接。

\teachervoice{讲者把这些方向明确列为 future work，包括去除 mini-batch approximation、few-shot generalization、domain adaptation、semi-supervised 与 multi-task learning。讲义只把它们作为研究路线，不把尚未完成的实验写成既成能力。}

\subsection{本章小结}

Dataset attention 的表达力与 $O(b^2)$ 代价来自同一设计。GNN、retrieval、sparse attention 与 meta-learning 提供自然扩展，但每种近似都会改变可见 context 与 statistical contract，需要新的机制验证。

